\section{The arithmetic boundary and the representation discipline}
\label{sec:foundations}

\subsection{The precise base theory}

Throughout, \iPRA{} means intuitionistic first-order arithmetic in the
language containing a symbol for every primitive recursive numerical
function, with their defining equations, successor injectivity and
separation, and induction for quantifier-free formulas. All variables in
an internal theorem range over natural numbers. The repository implements
exactly this theory as \texttt{PRAAx}; its primitive recursive symbols are
\texttt{PR k}, its arithmetic terms are \texttt{PTm n}, and its derivations
are \texttt{Prf PRAAx}. The native scalar presentation is connected to it by
\texttt{pra\_iff\_native} \cite{repo-pr,repo-native}.

An assertion about a primitive recursive presentation $A$ is a
\emph{schema}: fix its finite primitive recursive definition externally,
then substitute that definition into a first-order formula. Numerical
parameters of $A$, including degree, dimension, precision, and a finite
input polynomial, remain ordinary quantified variables. We will never
infer that a total computable function is primitive recursive, nor that a
primitive recursive predicate has its universal truth provable in
\iPRA. The latter would confuse representation with proof.

We use natural codes for signed integers, rational numbers with positive
denominator, finite tuples, and finite lists. One explicit list coding
suffices: let
\[
 \langle a,b\rangle=(a+b)^2+a+1,
 \qquad []=0,\qquad [a;L]=\langle a,L\rangle.
\]
Both components of a nonzero pair are less than its code. Recover them
by a search bounded by the code, with a default value on malformed
inputs. Injectivity follows by separating successive squares and then
subtracting. A traversal can be bounded by the initial code; an explicitly
stored length permits a sharper bound. All accesses, finite maxima,
appends, rational cross-multiplication tests, and their totalized default
branches are primitive recursive. These facts concern the numerical
algorithms; their algebraic laws follow by the corresponding bounded
inductions. Nothing depends on efficient binary encodings.

\begin{lemma}[Boolean reflection and bounded induction]
\label{lem:reflection}
For each bounded first-order formula $B(\mathbf x)$ built from primitive
recursive predicates there is a primitive recursive Boolean
$\chi_B(\mathbf x)\in\{0,1\}$ such that
\[
 \iPRA\vdash B(\mathbf x)\leftrightarrow \chi_B(\mathbf x)=1.
\]
Consequently induction on $B$ is derivable using quantifier-free induction.
The formula $B$ is fixed externally; this is not a uniform truth predicate
for all bounded formulas with arbitrary primitive recursive symbols.
\end{lemma}
\begin{proof}
Equality and numerical order have primitive recursive characteristic
functions. For example compare two numbers by simultaneous truncated
subtraction. Their defining recursions prove decidability by
quantifier-free induction. Use Boolean complement, conjunction, and
disjunction for the corresponding connectives. Since the already
reflected subformulas are decidable, implication is represented by
$\max(1-\chi_C,\chi_D)$; this use of Boolean implication makes no assertion
of excluded middle for arbitrary formulas.

A bounded universal quantifier is represented by the fold
\[
 F(0)=1,\qquad F(j+1)=F(j)\chi_C(j,\mathbf x).
\]
Here is a noncircular proof of its reflection law. With $i$ free,
quantifier-free induction on $j$ proves
$i<j\wedge F(j)=1\longrightarrow\chi_C(i,\mathbf x)=1$; the successor
step uses the cases $i<j$ and $i=j$. Conversely fix a final $n$ and
keep $\forall i<n\ \chi_C(i,\mathbf x)=1$ as a context assumption.
Induction on the quantifier-free formula $j\le n\longrightarrow F(j)=1$
proves $F(n)=1$. Thus bounded induction is not presupposed in this
proof of bounded reflection. The existential fold uses
$1-\prod(1-\chi_C)$. Its witness direction can also use the bounded
first-success scan described below. Formula induction here is an external
induction on a fixed syntax tree; only the folds' numerical inductions
occur inside \iPRA. Finally replace $B$ in its base and step by the
provably equivalent equation $\chi_B=1$, apply the permitted induction,
and replace back.
\end{proof}

\begin{lemma}[Bounded witness extraction]
\label{lem:bounded-choice}
Let $R(\mathbf x,i,y)$ be a fixed primitive recursive relation and
$b(\mathbf x,i)$ a fixed primitive recursive bound. There is a primitive
recursive $w(\mathbf x,i)$ such that \iPRA{} proves
\[
 (\exists y\le b(\mathbf x,i)\ R(\mathbf x,i,y))
 \ \longrightarrow\ R(\mathbf x,i,w(\mathbf x,i))
       \ \wedge\ w(\mathbf x,i)\le b(\mathbf x,i).
\]
For any numerical $m$, the list of these witnesses for $i<m$ is primitive
recursive and satisfies every required relation whenever the $m$ bounded
existences hold.
\end{lemma}
\begin{proof}
Scan $0,\ldots,b$, keeping a flag indicating whether a successful test
has occurred and, if so, its first index. The invariant says that either
no scanned test succeeds, or the saved index is in the scanned range,
satisfies $R$, and precedes every other success. All its quantifiers have
bounds supplied by the scan counter. Lemma~\ref{lem:reflection} permits
induction on this invariant. A supplied success forces the second case
at the end. In the no-success case return zero; this makes the operation
total. A second bounded fold appends the $m$ outputs. This is a construction
of a list, not an invocation of a choice axiom.
\end{proof}

\subsection{What a presented infinite object is}

A \emph{presentation} has three parts:
\begin{enumerate}
\item a fixed primitive recursive observation function
      $A(\mathbf a,j)$ with numerical parameters $\mathbf a$;
\item finite metadata, such as an amplitude bound, a dimension, a
      separation witness, or a primitive recursively supplied modulus;
\item arithmetic contracts on its observations. Frequently these have the
      form $\forall\mathbf u\ R(\mathbf a,\mathbf u)$ for a primitive
      recursive relation $R$.
\end{enumerate}
Examples include a finite matrix, an exact algebraic presentation, a
rational approximation to a complex number, the coefficients of a formal
power series, a finite probability law, and a rational sampling routine
for a function. A fixed function can have several parameters; that does
not turn it into a universal evaluator of arbitrary programs.

Different contracts express different mathematical content. A formal
power series needs no convergence modulus to state identities of finitely
many coefficients. Evaluation at a point additionally needs a quantitative
tail estimate. A Cauchy name permits approximate equality tests, but an
exact algebraic presentation can provide a decidable zero test. The
framework keeps these interfaces distinct because algorithms really do
use different information.

\begin{definition}[Regular rational names]
A fixed rational observation function $a(k)$ is regular when
\begin{equation}
 |a(k)-a(\ell)|\le 2^{-k}+2^{-\ell}
 \qquad(k,\ell\in\N).                         \tag{RN}\label{eq:RN}
\end{equation}
Complex names use pairs of rational observations and the rational
squared-norm version of the same estimate. We write $a\sim b$ when
\[
 \forall k\quad |a(k)-b(k)|\le 2^{1-k}.
\]
This is an arithmetic relation on two fixed presentations, not a
quotient type assumed as an object of first-order arithmetic. Its
usual real-number reading is explanatory notation.
\end{definition}

\begin{lemma}[Eliminating an arbitrary precision error]
\label{lem:vanishing-error}
For rational $q,c$ and nonnegative rational $C$, if
$q\le c+C2^{-r}$ for every $r$, then $q\le c$.
This implication is provable in \iPRA{} as a schema and requires neither
Markov's principle nor an unbounded minimization.
\end{lemma}
\begin{proof}
Compare the rational numbers $q,c$. If $q\le c$, finish. Otherwise write
$q-c=v/w$ with positive integers $v,w$, and write $C=a/b$ with $b\ge1$.
Choose, for example, $r=aw+2$. Then $2^r>aw$ (including $a=0$), so
$C2^{-r}<1/w\le v/w$. Substitution of this particular $r$ contradicts
the assumed inequality. The exponential inequality follows by ordinary
quantifier-free induction. The case distinction is rational and decidable.
\end{proof}

This lemma supplies elementary passages to exact name equality. For
example, transitivity of $\sim$ follows by comparing $a(k)$ and $c(k)$
through $a(r),b(r),c(r)$: regularity and the two equivalences give
\[
 |a(k)-c(k)|\le 2^{1-k}+6\,2^{-r}.
\]
Eliminate the last term by Lemma~\ref{lem:vanishing-error}. Thus no real
completion theorem is required for this equivalence law.

\begin{proposition}[Arithmetic on names with budgets]
\label{prop:name-arithmetic}
Regular names admit primitive recursive addition and multiplication with
explicit precision shifts. Inversion admits such a construction when a
positive apartness certificate is supplied.
\end{proposition}
\begin{proof}
For addition use $c(k)=a(k+2)+b(k+2)$ and add the two regularity bounds.
For multiplication suppose every sample of each input has absolute value
at most the supplied integer $M\ge1$. Choose $s$ with $2^s\ge2M$ and set
$c(k)=a(k+s)b(k+s)$. The identity
\[
 ab-a'b'=a(b-b')+b'(a-a')
\]
bounds the difference by $2M(2^{-k-s}+2^{-\ell-s})$, as required.
A cap can itself be derived from a zero-th sample and regularity:
$|a(k)|\le|a(0)|+2$.

For inversion suppose $|a(h+2)|\ge2^{-h}$. Every sample of index
$j\ge h+2$ then has $|a(j)|\ge2^{-h-1}$. Put $s=2h+3$ and
$c(k)=1/a(k+s)$, with a default zero value if a denominator test fails.
On valid inputs it does not fail, and
\[
 |c(k)-c(\ell)|\le
   2^{2h+2}(2^{-k-s}+2^{-\ell-s})
   \le2^{-k}+2^{-\ell}.
\]
The same finite difference identities show extensionality under $\sim$;
any remaining error is removed by Lemma~\ref{lem:vanishing-error}.
All regularity proofs fix the two requested precisions before any
induction on a finite sum or product.
\end{proof}

The apartness witness is substantive. The premise $\neg(a\sim0)$ does
not itself specify one. Likewise, a root-producing procedure may depend
on the presentation and choose different valid roots from equivalent
coefficient names. We demand extensionality of the asserted mathematical
relation, not an impossible global continuous choice of a preferred root.

\subsection{Three claims that must not be conflated}

We use the following status vocabulary throughout the report.
\begin{description}[style=nextline]
\item[Mathematical derivation in the stated arithmetic.]
A human-readable argument with primitive recursive constructions and
identified quantifier-free induction invariants. This is the main level
of the present contribution.
\item[Ambient Lean verification.]
A theorem checked by Lean about finite tables, streams, or arithmetic.
This checks that theorem in Lean's foundations. It does not by itself
produce a proof in the much weaker object theory.
\item[Deep iPRA verification.]
An actual term of type $\Gamma\mathbin{\vdash}[\mathtt{PRAAx}]\varphi$,
with the primitive recursive definitions and the formula $\varphi$
explicitly encoded. The repository already provides the logical backend;
the new mathematical compilation layer is not yet fully formalized.
\end{description}

The distinction is constructive, not merely terminological: a Lean proof
by induction on an arbitrary proposition does not authorize that induction
inside \iPRA. The next section gives the finite invariant that makes
specific uses legitimate.
